% BEGIN PROFESSOR FEEDBACK PF15: Implementation - Project structure
\subsection{Project structure and entry points}

The repository keeps the exploratory and the evidence-producing code apart. The
interactive command-line interface (CLI) in \texttt{main.py} is for development
and inspection, while the thesis
results come only from \texttt{scripts/run\_results.py}, whose \texttt{prepare} and
\texttt{execute} subcommands call the separate \texttt{src/experiments} package.
The same entry point provides read-only \texttt{preflight} and \texttt{verify}
commands, execution \texttt{--resume}, and a coordinated
\texttt{execute-parallel} command for running the HF and Real-CATS roots on one
server.
Keeping the two apart stops exploratory shortcuts, such as row limits or changed
defaults, from leaking into the results. The main files are shown below, and
Figure~\ref{fig:pipeline_architecture} shows how they connect across the two
phases of a run.

\begin{lstlisting}[caption={Project layout (abbreviated).},
  label={lst:layout}]
bachelor_thesis/
main.py              CLI entry point, pipeline orchestrator
config.py            Exploratory defaults only
requirements.txt     Environment definition
requirements-results.txt Exact numerical Results environment
scripts/run_results.py  Locked Results prepare/execute CLI
src/
 experiments/
   constants.py     Locked constants and verified role catalogue
   prepare.py       Source conversion, cohorts, folds, protocol lock
   datasets.py      HuggingFace streaming and BLTE loading
   realcats.py      SQLite deduplication and layer export
   ordering.py      DuckDB external event ordering
   features.py      Baseline and fingerprint transformations
   motifs.py        Arrow-to-Raphtory loading and motif sweep
   modeling.py      Fixed estimators and development selection
   statistics.py    Metrics, permutations, bootstrap, BH correction
   splits.py        Group-aware and stratified partition construction
   rq_binary.py     RQ1/RQ2 execution and persistence
   rq_multiclass.py RQ3/RQ4 execution and persistence
   secondary.py     Conditional BLTE and synthetic validation
   synthetic_fixtures.py  Local-role and boundary validation templates
   execute.py       Verification and RQ orchestration
   execution_control.py Locks, attempts, threads, state transitions
   execution_artifacts.py Checkpoints, stages, result inventory
   coordination.py  Same-host HF/Real-CATS coordination
   verify.py        Read-only integrity and citability checks
   validation.py    Prepared-data and result-schema guards
   artifacts.py     Exclusive writes, hashes, runtime provenance
   resources.py     Read-only memory/disk projection and admission
   telemetry.py     Timing log and five-minute heartbeat
   types.py         Shared data contracts (dataclasses)
   visualization.py Chapter-ready figures
 graph_construction.py  Exploratory NetworkX graph features
 motifs/            Exploratory alternative motif backends
data/raw/            Raw datasets (gitignored)
results/experiments/ Locked run directories
\end{lstlisting}

\begin{figure}[p]
  \centering
  \ResGraphic{1.0}{01-pipeline}
  \caption{The locked Results pipeline. The preparation phase fixes the cohorts
  and writes the locked artifacts, and the execution phase then computes the
  local role counts, trains the models, and runs the statistical tests. Italic labels name
  the source modules in \texttt{src/experiments}.}
  \label{fig:pipeline_architecture}
\end{figure}
% END PROFESSOR FEEDBACK PF15

